\frfilename{mt13.tex}
\versiondate{16.6.01}
\copyrightdate{1995}

\def\chaptername{Pelengkap}
\def\sectionname{Pendahuluan}

\newchapter{13}

Dalam bab ini saya menghimpun sejumlah hasil yang tidak terletak pada alur
langsung argumen dari 111A (definisi `aljabar-$\sigma$') menuju 123C
(Teorema Konvergensi Terdominasi Lebesgue), tetapi tetap harus disertakan
dalam jilid ini karena semuanya relatif elementer, esensial bagi pengembangan
lebih lanjut, dan diperlukan untuk memahami secara tepat apa yang telah
dikerjakan. Bagian terpanjang adalah \S134, yang membahas beberapa sifat
khusus elementer ukuran Lebesgue; khususnya, invariansinya terhadap translasi,
keberadaan himpunan dan fungsi tak-terukur, serta himpunan Cantor. Bagian-bagian
lainnya relatif ringan. \S131 membahas subruang (terukur) dan penafsiran
rumus $\int_Ef$, yang memperumum gagasan integral $\int_a^bf$ suatu fungsi
pada interval. \S132 memperkenalkan ukuran luar yang terkait dengan suatu
ukuran, semacam kebalikan dari konstruksi \Caratheodory{} untuk membangun ukuran
dari ukuran luar. \S\S133 dan 135 menetapkan konvensi yang sesuai untuk menangani
`tak hingga' dan bilangan kompleks (secara terpisah! keduanya tidak berpadu
dengan baik) sebagai nilai integran ataupun integral; sekaligus saya
menyebutkan integral `atas' dan `bawah'. Terakhir, dalam \S136, saya memberikan
beberapa teorema mengenai aljabar-$\sigma$ himpunan, yang menguraikan bagaimana
aljabar tersebut (dalam beberapa kasus terpenting) dapat dibangkitkan melalui
operasi-operasi yang relatif terbatas.

\discrpage
